\documentclass[aspectratio=169,10pt]{beamer}

\usepackage[T1]{fontenc}
\usepackage{lmodern}
\usepackage{amsmath,amssymb}
\usepackage{hyperref}

% ---------------------------------------------------------------- design
\definecolor{kred}{RGB}{178,28,36}
\definecolor{kgrey}{RGB}{90,90,90}

\usetheme{default}
\usecolortheme{default}
\setbeamercolor{structure}{fg=kred}
\setbeamercolor{frametitle}{fg=kred,bg=white}
\setbeamercolor{title}{fg=kred}
\setbeamercolor{itemize item}{fg=kred}
\setbeamertemplate{navigation symbols}{}
\setbeamertemplate{itemize item}{\raisebox{0.15ex}{\tiny$\blacksquare$}}
\setbeamerfont{frametitle}{series=\bfseries,size=\large}
\setbeamerfont{title}{series=\bfseries}
\setbeamerfont{itemize/enumerate body}{size=\small}
\setbeamertemplate{footline}{%
  \hspace{2ex}\textcolor{kgrey}{\tiny Halil \.{I}brahim Kanpak, 27 October 2026}%
  \hfill\textcolor{kgrey}{\tiny\insertframenumber}\hspace{2ex}\vspace{1ex}}
\setlength{\leftmargini}{1.4em}

\hypersetup{colorlinks=true,urlcolor=kred,linkcolor=kred}

\newcommand{\redline}{\textcolor{kred}{\rule{\textwidth}{1pt}}}
% one work per frame: venue and authors, link, setting, bound
\newcommand{\venue}[3]{\item #1. #2.\\ \href{#3}{\nolinkurl{#3}}}

\title{meeting 2026-10-27}
\author{Halil \.{I}brahim Kanpak}
\date{2026-10-27}

\begin{document}

% ================================================================ title
\begin{frame}[plain]
  \vfill
  {\Huge\textbf{\textcolor{kred}{meeting 2026-10-27}}}\\[0.4em]
  {\large Unbounded and polynomial-time attackers in learning security and privacy}\\[1.2em]
  \redline\\[0.8em]
  {\small Halil \.{I}brahim Kanpak, Ko\c{c} University}
  \vfill
\end{frame}

% ================================================================ terms
\begin{frame}{Two attackers and differential privacy}
  \begin{itemize}
    \item An unbounded attacker has no limit on its running time. A polynomial-time attacker runs in time polynomial in the security parameter.
    \item A curator holds a database and answers queries through a mechanism $M$, a randomized algorithm. Two databases are neighbours if they differ in one record. $M$ is $(\varepsilon,\delta)$-differentially private (DP) if $\Pr[M(D)\in S]\le e^{\varepsilon}\Pr[M(D')\in S]+\delta$ for all neighbours $D,D'$ and every set $S$ of outputs.
    \item Statistical DP (SDP) is this definition. It holds against an unbounded attacker, since it covers every set of outputs. Computational DP (CDP) asks the same inequality only for the tests a polynomial-time attacker can run (Mironov et al., CRYPTO 2009).
  \end{itemize}
\end{frame}

% ================================================================ state
\begin{frame}{The state of the literature}
  \begin{itemize}
    \item A task is a query with a utility function that scores answers. A task separates CDP from SDP when some CDP mechanism has good utility on it and no SDP mechanism does, efficient or not. Such a separation is proven when two parties each hold part of the data. With a single curator the only separation we know is a task built for that purpose (proven, under strong cryptographic assumptions).
    \item In robustness, backdoors and watermarking, each result we found where the two attackers differ hides a cryptographic object in the data or the model. A digital signature is the typical example (empirical, on our searches).
    \item For membership inference we found no hardness result that rests on a cryptographic assumption. For recovering the input of a network with known weights, only worst-case NP-hardness is known (empirical, on our searches).
  \end{itemize}
  \vspace{0.2em}
  \redline\\[0.2em]
  {\small Today we go through eleven works. For each we give the setting, and say which party the authors limit in running time and how.}
\end{frame}

% ================================================================ 1
\begin{frame}{Computational Differential Privacy}
  \begin{itemize}
    \venue{CRYPTO 2009}{Mironov, Pandey, Reingold, Vadhan}{https://doi.org/10.1007/978-3-642-03356-8_8}
    \item A curator holds a database in which each record belongs to one person. A mechanism answers a query with a randomized output, and the attacker sees the output. The paper also studies a two-party setting. Alice and Bob each hold a bit vector and run a protocol to estimate the Hamming distance, the number of positions where the two vectors differ. A party's view is everything it sees during the protocol. Changing one bit of Alice's vector may change the distribution of Bob's view only by the factor DP allows, and the same holds with the roles swapped.
    \item SDP asks DP against every attacker, whatever its running time. The paper defines CDP against non-uniform polynomial-time attackers, in two forms. IND-CDP asks the DP inequality only for the tests such an attacker can run, up to a negligible term. SIM-CDP asks that the mechanism's output be computationally indistinguishable from the output of some SDP mechanism on the same database. SIM-CDP implies IND-CDP (proven). For the Hamming distance the paper's SDP protocol has an error that grows with the distance, and its SIM-CDP protocols have an error independent of the vector length (proven).
  \end{itemize}
\end{frame}

% ================================================================ 2
\begin{frame}{The Limits of Two-Party Differential Privacy}
  \begin{itemize}
    \venue{FOCS 2010, corrected version ECCC TR11-106}{McGregor, Mironov, Pitassi, Reingold, Talwar, Vadhan}{https://eccc.weizmann.ac.il/report/2011/106/}
    \item Alice and Bob each hold an $n$-bit vector. They run a protocol to estimate the Hamming distance or the inner product of the two vectors. Changing one bit of one party's vector may change the distribution of the other party's view only by the factor DP allows, so each party is the attacker against the other. With a single curator who holds both vectors, the Hamming distance can be released with constant error.
    \item Against unbounded attackers, every such two-party protocol has additive error of order $\sqrt{n}$, up to a logarithmic factor, and this is tight (proven). The bound holds for $\varepsilon$-DP and for $(\varepsilon,\delta)$-DP with $\delta=o(1/n)$. Against polynomial-time attackers, the parties can use secure function evaluation, a cryptographic protocol that computes a function of both inputs and reveals only its output. With it they simulate a trusted curator who adds Laplace noise, so the error under CDP is constant in $n$ (proven by Mironov et al.). This separates SDP from CDP when two parties hold the data.
  \end{itemize}
\end{frame}

% ================================================================ 3
\begin{frame}{Limits of Computational Differential Privacy in the Client/Server Setting}
  \begin{itemize}
    \venue{TCC 2011}{Groce, Katz, Yerukhimovich}{https://doi.org/10.1007/978-3-642-19571-6_25}
    \item A server holds the whole database. A client sends queries and receives the answers, and the client is the attacker. A task is a query together with a utility function, which scores how useful an answer is for the database.
    \item The paper asks whether a polynomial-time attacker lets the server do better than an unbounded one, with more useful answers or with faster mechanisms. It gives two negative answers. Take a query whose answer is a vector of constant dimension, and a utility that measures the error of the answer in an $\ell_p$ norm or a generalization of it. Then any CDP mechanism can be turned into an SDP mechanism with almost the same utility and running time (proven). Take instead a CDP mechanism that uses trapdoor permutations, collision-resistant hash functions or random oracles only as black boxes. Then an SDP mechanism exists that is as useful and as fast, for any utility and any output (proven).
  \end{itemize}
\end{frame}

% ================================================================ 4
\begin{frame}{Towards Separating Computational and Statistical Differential Privacy}
  \begin{itemize}
    \venue{FOCS 2023}{Ghazi, Ilango, Kamath, Kumar, Manurangsi}{https://arxiv.org/abs/2301.00104}
    \item One curator holds a database of $n$ bits, and two databases are neighbours if they differ in one bit. The curator must output a circuit. The circuit is useful if it accepts the database and rejects every string in a set $R$ that is at Hamming distance more than $\tau$ from the database. A keyless hash function is one fixed public hash function, and $R$ holds the strings it maps to zero. The output is a circuit, so it is high-dimensional, and the constant-dimension result of Groce et al. does not apply.
    \item Some mechanism is $\varepsilon$-IND-CDP with constant $\varepsilon$ and is useful with probability $1-o(1)$. Every $(\varepsilon,\delta)$-SDP mechanism that is useful with constant probability, even an inefficient one, needs $\varepsilon$ to grow with $n$ or $\delta$ to be at least an inverse polynomial (proven under three assumptions). The first is non-interactive witness-indistinguishable proofs. The second is that the keyless hash has short output and that an efficient attacker finds only polynomially many collisions. The third is differing-inputs obfuscation, which makes two circuits look the same unless an attacker can find an input on which they differ. The proof first shows privacy against attackers that query the circuit at points fixed in advance. The keyless hash extends it to attackers that choose their queries adaptively, and obfuscation extends it to every polynomial-time attacker that sees the circuit.
  \end{itemize}
\end{frame}

% ================================================================ 5
\begin{frame}{Adversarially Robust Learning Could Leverage Computational Hardness}
  \begin{itemize}
    \venue{ALT 2020}{Garg, Jha, Mahloujifar, Mahmoody}{https://arxiv.org/abs/1905.11564}
    \item A learner trains a classifier on samples from a distribution, and a challenger draws a fresh test instance. The attacker gets the instance and query access to the classifier, and may flip a few bits of the instance. It wins if the classifier gives the changed instance a wrong label. In the first construction the classifier may also reject an input as tampered, and a rejected input is not a win.
    \item Robustness here is asked only against polynomial-time attackers. Each instance carries a short digital signature, a tag that anyone can verify and only the key holder can produce. The classifier rejects an instance whose signature is invalid, so the attacker needs a valid signature for a new instance. An unbounded attacker can forge one and a polynomial-time attacker cannot (proven, from exponentially hard one-way functions). A second construction has no rejection. Its error under attack is negligible against polynomial-time attackers and at least $1/2$ against unbounded attackers who flip even fewer bits (proven). Conversely, any such gap on a task with efficiently samplable instances gives hard-on-average problems in NP (proven).
  \end{itemize}
\end{frame}

% ================================================================ 6
\begin{frame}{Planting Undetectable Backdoors in Machine Learning Models}
  \begin{itemize}
    \venue{FOCS 2022}{Goldwasser, Kim, Vaikuntanathan, Zamir}{https://arxiv.org/abs/2204.06974}
    \item A user hands training to an untrusted trainer, who returns a classifier. A malicious trainer can plant a backdoor, a hidden behaviour tied to a secret backdoor key. With the key, the trainer can change the label of any input by changing the input only slightly. The user is the attacker against undetectability and tries to tell a backdoored classifier from an honest one. In the black-box case the user only queries the classifier. In the white-box case the user also sees the weights and the training data.
    \item Undetectability is asked only against polynomial-time attackers. The black-box backdoor works for any classifier and uses a digital signature, so only inputs that carry a valid signature trigger it. Finding any input where the backdoored and honest classifiers differ is then infeasible (proven, from one-way functions). The white-box backdoor is for random Fourier features, a model whose first layer has random weights. The trainer only changes how these weights are drawn, and detecting the change would break continuous LWE, a variant of LWE with real-valued samples (proven). A second white-box variant, for random ReLU networks, rests on sparse PCA, the problem of finding a hidden sparse direction in noisy samples. Our round of 26 October found it detectable at the printed parameters (proven by us, not yet checked by a person).
  \end{itemize}
\end{frame}

% ================================================================ 7
\begin{frame}{Statistically Undetectable Backdoors in Deep Neural Networks}
  \begin{itemize}
    \venue{ICML 2026}{Bogdanov, Rosen, Vafa}{https://proceedings.mlr.press/v306/bogdanov26a.html}
    \item A trainer releases a feedforward network on integer inputs. Its first layer multiplies the input by a random Gaussian matrix with fewer rows than columns. The later layers are bilipschitz, which means they shrink or stretch distances only by factors between two fixed bounds. The backdoored network draws this matrix together with a secret vector of entries $+1$ and $-1$ that the matrix sends close to zero. This vector is the backdoor key. Adding the key to an input moves the input far but barely changes the output. The user sees every weight and is the attacker.
    \item The backdoored weights and the honest weights have almost the same distribution, at total variation distance $o(1)$ (proven, with no assumption). So no attacker, whatever its running time, detects the backdoor with more than $o(1)$ advantage. Without the key, a polynomial-time attacker cannot find two distant inputs with nearly equal outputs. This rests on the hardness of finding short integer vectors that a Gaussian matrix sends close to zero, which the authors tie to worst-case lattice problems (proven under that assumption).
  \end{itemize}
\end{frame}

% ================================================================ 8
\begin{frame}{Backdoor Defense, Learnability and Obfuscation}
  \begin{itemize}
    \venue{ITCS 2025}{Christiano, Hilton, Lecomte, Xu}{https://arxiv.org/abs/2409.03077}
    \item The paper defines backdoor defense as a game between an attacker and a defender. The attacker picks a distribution over inputs and a function $f$ from a class $F$. A trigger $x^*$ is then drawn at random from the distribution. The attacker returns a backdoored function $f^*$ in $F$ that differs from $f$ at $x^*$, and on at most an $\varepsilon$ fraction of inputs drawn from the distribution. The defender gets a function and an input. It must decide whether it holds the clean pair $(f,x)$ with a random $x$ or the backdoored pair $(f^*,x^*)$, and it wins if it decides correctly.
    \item Here the limit on running time sits on the defender. With no limit, a defender wins with probability tending to 1 exactly when $\varepsilon=o(1/\mathrm{VC}(F))$ (proven). The Vapnik-Chervonenkis dimension $\mathrm{VC}(F)$ is the size of the largest set of inputs on which $F$ realizes every labelling. It also decides PAC learnability, so defense with no time limit behaves like learning. A polynomial-time defender wins for every class that is efficiently PAC learnable (proven). In the random oracle model it also wins for some classes that are not efficiently learnable (proven). It loses on polynomial-size circuits if one-way functions and indistinguishability obfuscation exist (proven under these assumptions).
  \end{itemize}
\end{frame}

% ================================================================ 9
\begin{frame}{Undetectable Watermarks for Language Models}
  \begin{itemize}
    \venue{COLT 2024}{Christ, Gunn, Zamir}{https://arxiv.org/abs/2306.09194}
    \item A provider runs a language model and embeds a watermark in the generated text with a secret key. The user is the attacker. It sends prompts of its choice, one after another, and tries to tell whether it talks to the watermarked model or to the original model. The key holder detects the watermark in any long enough piece of text that was generated with enough randomness. The scheme is sound, which means text written without the key is flagged only with negligible probability.
    \item Against polynomial-time attackers the watermark is undetectable if one-way functions exist (proven). Against an unbounded attacker no watermark is. Take any sound watermark that the key holder detects with non-negligible probability on some prompt. An attacker with polynomially many queries and time exponential in the key length tells the watermarked model from the original (proven, Lemma 4). So undetectability can hold only against attackers with bounded running time.
  \end{itemize}
\end{frame}

% ================================================================ 10
\begin{frame}{Cryptanalytic Extraction of Neural Network Models}
  \begin{itemize}
    \venue{CRYPTO 2020}{Carlini, Jagielski, Mironov}{https://arxiv.org/abs/2003.04884}
    \item A provider exposes a ReLU network as an oracle. The attacker chooses inputs and receives the real-valued outputs, with no timing or other side channel. The attacker knows the number of layers and their widths. It wants a copy whose outputs match the oracle on all inputs up to a small error. This task is model extraction.
    \item The paper counts the attacker's queries and its running time separately. A ReLU network is piecewise linear. At a critical point one neuron switches between on and off, and the outputs near that point reveal the neuron's weights up to a scale factor whose sign is unknown. Recovering the signs of one layer takes polynomially many queries but time exponential in the layer width, because the attack tries every sign pattern of the layer (heuristic). No lower bound is known on the cost of extraction (open). The attack copied a 100,000-parameter MNIST network with $2^{21.5}$ queries, with worst-case output error $2^{-25}$ (empirical).
  \end{itemize}
\end{frame}

% ================================================================ 11
\begin{frame}{Polynomial Time Cryptanalytic Extraction of Neural Network Models}
  \begin{itemize}
    \venue{EUROCRYPT 2024}{Canales-Mart\'{\i}nez, Ch\'avez-Saab, Hambitzer, Rodr\'{\i}guez-Henr\'{\i}quez, Satpute, Shamir}{https://doi.org/10.1007/978-3-031-58734-4_1}
    \item The setting is the one of Carlini et al. The network is fully connected with ReLU activations, and the attacker knows the layer widths. The attacker queries chosen real-valued inputs and sees the logits, which are the outputs before softmax, or the softmax probabilities. The analysis assumes arithmetic with infinite precision.
    \item The paper removes the exponential search over signs. Its main new attack, neuron wiggling, starts at a critical point of one neuron and moves the input in the two opposite directions that change that neuron most. The ReLU blocks the change on one side only, so the side with the larger output change reveals the sign. The attack repeats this at many inputs, about 200 in the experiments, and decides by a statistical vote. Two other techniques handle the first two hidden layers and the last one. Extraction then takes polynomially many queries and polynomial time (heuristic, because the authors exclude rare failures such as neurons that almost never switch). On a CIFAR-10 network with 8 hidden layers of 256 neurons, sign recovery took about 30 minutes on a 256-core computer (empirical).
  \end{itemize}
\end{frame}

\end{document}
